\chapter{Matrix calculus and Jacobians}
\label{ch:calc}

Generative models differentiate scalars with respect to vectors and matrices. This chapter fixes the conventions and derives the four derivatives the homeworks need: a linear form, a quadratic form, a log-determinant, and a Jacobian determinant.

\section{Layout}

A scalar $L(\theta)$ has a gradient $\nabla_\theta L$ that is a column vector of the same shape as $\theta$. The differential is the inner product
\[
dL=\langle \nabla_\theta L,\, d\theta\rangle.
\]
If $\theta$ is a matrix, the inner product is the Frobenius product $\langle A,B\rangle=\tr(A^\top B)$. This is the layout used below. It matches the chain rule in coordinates.

\begin{resultbox}[Linear and quadratic forms]
Let $a\in\R^d$ be constant and let $A\in\R^{d\times d}$ be constant. Then
\begin{align*}
\nabla_x (a^\top x)&=a,\\
\nabla_x (x^\top A x)&=(A+A^\top)x.
\end{align*}
If $A$ is symmetric, the second gradient is $2Ax$. In particular $\nabla_x \|x\|^2=2x$ and $\nabla_x \tfrac12\|x\|^2=x$.
\end{resultbox}

\begin{intbox}
The factor $2$ is the usual one from $(x^2)'=2x$. People drop it when they differentiate $\tfrac12 x^\top \Sigma^{-1}x$, because the $\tfrac12$ cancels the $2$. Gaussian exponents are written with that $\tfrac12$ for this reason.
\end{intbox}

\section{Trace and log-determinant}

\begin{resultbox}[Derivatives in the matrix]
For constant $B$ and invertible $A$,
\begin{align*}
\nabla_A \tr(A^\top B)&=B,\\
\nabla_A \log\lvert\det A\rvert&=A^{-\top}.
\end{align*}
The second identity is the matrix calculus behind Glow. If $W$ depends on a diagonal vector $s$ only through $W=\diag(s)$, then
\[
\frac{\partial}{\partial s_i}\log\lvert\det W\rvert=\frac{1}{s_i}
\quad\text{when }s_i\neq 0,
\]
and the gradient of $\sum_i\log\lvert s_i\rvert$ is the vector $(1/s_i)$.
\end{resultbox}

A short derivation of the log-determinant: $\det(A+dA)=\det(A)\det(I+A^{-1}dA)$, and $\det(I+H)=1+\tr(H)+o(\|H\|)$. Therefore
\[
d\log\lvert\det A\rvert=\tr(A^{-1}\,dA)=\langle A^{-\top},\, dA\rangle.
\]

\section{Jacobians and the change-of-variables factor}

Let $f:\R^D\to\R^D$ be continuously differentiable and invertible, with differentiable inverse, and with $J_f(u)$ invertible at the points we use. If $u$ has density $p_0$ and $y=f(u)$, then $y$ has density
\begin{equation}
\label{eq:cov}
p(y)=p_0(f^{-1}(y))\,\lvert\det J_{f^{-1}}(y)\rvert
=p_0(u)\,\lvert\det J_f(u)\rvert^{-1}.
\end{equation}
In logarithms, with $u=f^{-1}(y)$,
\[
\log p(y)=\log p_0(u)-\log\lvert\det J_f(u)\rvert.
\]

\begin{intbox}[Why the absolute value]
A reflection $f(u)=-u$ preserves probability but has Jacobian $-1$. Densities are nonnegative, so the volume factor is the absolute value of the determinant. On a triangular matrix the determinant is the product of the diagonal entries, which is why coupling layers are easy.
\end{intbox}

\begin{resultbox}[Triangular determinant]
If $L$ is lower triangular with ones on the diagonal, $\det L=1$. If $U$ is strictly upper triangular, then $U+\diag(s)$ is upper triangular with diagonal $s$, and $\det(U+\diag(s))=\prod_i s_i$. A permutation matrix $P$ has $\det P=\pm 1$. Hence if $W=PL(U+\diag(s))$,
\[
\log\lvert\det W\rvert=\sum_{i=1}^{C}\log\lvert s_i\rvert.
\]
\end{resultbox}

\begin{practbox}[A $2\times 2$ Jacobian]
Let $f(u_1,u_2)=(u_1,\, u_2+u_1^2)$. Compute $J_f$, its determinant, and the density of $y=f(u)$ when $u\sim\N(0,I_2)$.
\end{practbox}
\begin{solbox}
\[
J_f(u)=\begin{pmatrix}1&0\\ 2u_1&1\end{pmatrix},\qquad \det J_f=1.
\]
The map is a shear. It preserves volume, so $\log p(y)=\log p_0(f^{-1}(y))$. The inverse is $(y_1,\, y_2-y_1^2)$, and
\[
p(y)=(2\pi)^{-1}\exp\bigl(-\tfrac12 y_1^2-\tfrac12(y_2-y_1^2)^2\bigr).
\]
This is the map in Homework~2, problem~3.2. It will return there as a counterexample to a diagonal-Gaussian restriction.
\end{solbox}
